\begin{textsample}{NEG Dim 1 – global\_south\_2025\_06 – Score 1.00 – global\_south\_2025\_06/126301780.txt}  \label{ex:f1_neg_239}
The West remains silent on Gaza ; we mustn't Dhaka -- I ' ve been trying to follow the Global March to Gaza for over a month now.
The only glimpses I ' ve found come from Al Jazeera, Turkiye ' s Anadolu Ajansi, and a handful of Palestinian news portals.
But the march is happening -- right now -- and most of the Western world either doesn't seem to notice or blatantly choosing not to.
Across every inhabited continent, thousands are mobilizing : marching, chanting, risking arrest, demanding justice for Gaza.
Among them are doctors and teachers, imams and priests, students and elders -- united in historic solidarity.
And yet, if you turned on your television this morning, scrolled through your feed, or flipped through a newspaper, you ' d barely know it was happening.
This silence is not accidental.
It is carefully engineered.
Article continues after this advertisement Picture this : a convoy of nine buses carrying nearly a thousand activists, backed by @ @ @ @ @ @ @ @ @ @ Tunis and crossing into Libya -- en route to Egypt ' s Rafah crossing.
They ' re supported by Tunisian labor leaders and legal defenders, standing in defiance of a siege in what the United Nations calls"the hungriest place on earth."Their goal is symbolic and profound : to break the blockade, deliver aid, bear witness to genocide, and pressure world leaders into action -- much like the Madleen crew, whose mission was intercepted just days ago.
This is a truly global movement.
Activists from 50 countries are converging in Cairo under banners like the Palestinian Youth Movement, Codepink, and Jewish Voice for Labor -- with support from 150 organizations across 31 nations.
Hundreds more are joining from Algeria, Libya, Morocco, and Mauritania.
Some will come by land, others by sea.
The stakes could not be higher.
Gaza has been under siege for more than 20 months.
A genocidal war has killed over 54,900 Palestinians and injured more than 126,000.
Aid distribution sites have become @ @ @ @ @ @ @ @ @ @ gathering food since May 27, wounding nearly 2,000.
Schools lie in ruins, hospitals are bombed out, water is scarce.
Hunger is weaponized.
Gaza is collapsing.
Article continues after this advertisement And still, the Western media says almost nothing.
This diverse, global act of civilian resistance doesn't fit the dominant narrative where Palestine is framed as fringe insurgency, Israel as self-defense.
To acknowledge thousands marching peacefully against genocide would dismantle that fiction.
So, the silence holds -- no live feeds, no front pages, no analysis.
Just vacuum.
A sickening and deliberate vacuum.
In fact, over 100 BBC employees have allegedly accused the British broadcaster of pro-Israel bias in its coverage of the Gaza war.
The claim was made in an open letter sent to BBC director Tim Davie and signed by more than 230 figures in the UK ' s media industry and other sectors, who said the public broadcaster has failed to provide"fair and accurate"coverage of the conflict. @ @ @ @ @ @ @ @ @ @ the United States -- the world ' s most powerful enabler of this violence.
Washington sends \$3.8 billion in unconditioned military aid to Israel each year.
Its vetoes at the UN Security Council blocks ceasefires and investigations.
Its diplomatic shield allows siege, bombing, displacement, and starvation to continue with impunity.
And because the US sits at the center of global media power, it shapes narratives at will.
Palestinian suffering is erased.
Israeli violence is sanitized.
Headlines mimic government talking points, not witness testimony.
What you see is not journalism ; it ' s narrative warfare.
But silence can not contain solidarity.
People are marching and acting how they can.
And you can too, in your own way.
Share livestreams ; amplify images ; write to political leaders and demand an end to arms transfers, call for ceasefire and UN investigations, insist on open humanitarian corridors ; organize vigils, teach-ins, divestment campaigns, legal actions ; support grassroots Palestinian media and aid organizations.
Recognize your abilities, your @ @ @ @ @ @ @ @ @ @ all have roles to play -- and every single act matters.
Because the choice is clear : in the face of genocide and silence, neutrality is not peace.
It is violence.
Standing by means siding with power.
But going over your own head will also be of no help.
Extend help within your means -- but extend it anyhow.
All of this combined has now become a moral reckoning.
When people touch the cables of silence, they spark cracks.
When voices break through the cover-up, truth begins to seep.
If the world ' s media won't show the march, we must.
If the powerful won't act, we must.
If neutrality is not an option, then action becomes an obligation.
The march is underway.
So is history.
The Daily Star/Asia News Network

% matched lemmas: 
\end{textsample}
